% 第0章: 前提知識 — C++クラスの基礎
\section{前提知識 — C++クラスの基礎}
\label{sec:prerequisite}

この章では、本書を読むために必要なC++クラスの基礎知識を解説します。

\begin{notebox}[すでにクラスを理解している方へ]
構造体・クラス・コンストラクタ・アクセス修飾子を理解している方は、この章を飛ばして次の「全部main.cppに書くとどうなる？」に進んでください。
\end{notebox}

%% ── Q0: プログラミングの部品を整理しよう ──────
\begin{qabox}{プログラミングの部品を整理しよう}

C++にはいろいろな「部品」が出てきますが、実はどれも「何かを入れる箱」としてイメージできます。
まずは全体像をつかみましょう。

\subsection{変数 — 値の箱}

一番シンプルな部品です。1つの値を入れておく箱。

\begin{lstlisting}
int speed = 100;       // 整数を1つ入れる箱
float temperature = 25.3; // 小数を1つ入れる箱
bool isRunning = true; // true/falseを1つ入れる箱
\end{lstlisting}

\subsection{配列 — 箱の列}

同じ種類の箱を一列に並べたもの。番号（インデックス）で「何番目の箱か」を指定してアクセスします。

\begin{lstlisting}
int scores[3] = {80, 95, 70};  // int型の箱が3つ並んだ列
scores[0];  // → 80（0番目）
scores[1];  // → 95（1番目）
\end{lstlisting}

\subsection{関数 — 処理の箱}

「一連の処理」をまとめて名前を付けた箱です。呼び出すと中の処理が実行されます。

\begin{lstlisting}
// 関数の定義（処理を箱に入れる）
void blinkLed(int pin) {
    digitalWrite(pin, HIGH);
    delay(500);
    digitalWrite(pin, LOW);
}

// 関数の呼び出し（箱を使う）
blinkLed(13);  // ピン13でLEDを点滅
\end{lstlisting}

\subsection{構造体 — 変数の箱}

関連する変数をまとめて1つの「箱」にしたもの。中に複数の変数を持てます。

\begin{lstlisting}
struct LedData {           // 「LedDataという箱の設計図」
    bool state = false;    //   中に入る変数1
    uint8_t brightness = 0; //   中に入る変数2
};
\end{lstlisting}

\subsection{クラス — 変数と関数の箱}

構造体の進化版。変数（データ）だけでなく、関数（操作）も一緒にまとめた箱です。

\begin{lstlisting}
class LedModule {          // 「LedModuleという箱の設計図」
    uint8_t _pin;          //   中に入る変数（データ）
    bool _state;           //   中に入る変数（データ）
    void turnOn();         //   中に入る関数（操作）
    void turnOff();        //   中に入る関数（操作）
};
\end{lstlisting}

\vspace{2mm}
ここまでの関係をまとめると、こうなります。

\begin{center}
\begin{tikzpicture}[
    box/.style={draw, rounded corners=3pt, minimum width=28mm, minimum height=10mm, align=center, font=\small\bfseries},
    arrow/.style={-{Stealth[length=2.5mm]}, thick, color=accentblue},
    label/.style={font=\scriptsize\color{gray}, align=center},
]
    \node[box, fill=accentblue!10] (var) at (0,0) {変数};
    \node[box, fill=accentblue!10] (arr) at (3.5,0) {配列};
    \node[box, fill=accentblue!10] (func) at (7,0) {関数};
    \node[box, fill=okgreen!15] (struct) at (2,-2) {構造体};
    \node[box, fill=okgreen!20] (class) at (6,-2) {クラス};

    \node[label, above=1mm of var] {値の箱};
    \node[label, above=1mm of arr] {箱の列};
    \node[label, above=1mm of func] {処理の箱};
    \node[label, below=1mm of struct] {変数の箱};
    \node[label, below=1mm of class] {変数＋関数の箱};

    \draw[arrow] (var) -- node[label, right, pos=0.5] {まとめる} (struct);
    \draw[arrow] (struct) -- node[label, above, pos=0.5] {関数も追加} (class);
    \draw[arrow] (func) -- node[label, right, pos=0.5] {組み込む} (class);
\end{tikzpicture}
\end{center}

この先の章では、構造体とクラスを詳しく見ていきます。

\end{qabox}

%% ── Q1: 構造体 ──────────────────────────────
\begin{qabox}{構造体（struct）って何？}

プログラムを書いていると、「関連する変数をセットで扱いたい」場面がよくあります。

たとえば、LEDの状態を管理する変数が2つあるとします。

\begin{badexample}{バラバラの変数で管理}
\begin{lstlisting}
bool ledState = false;
uint8_t ledBrightness = 0;

bool led2State = false;
uint8_t led2Brightness = 0;
// LED が増えるたびに変数が倍増...
\end{lstlisting}
\end{badexample}

LEDが増えるたびに変数名を考えなくてはいけませんし、「どの変数がどのLEDのものか」が分かりにくくなります。

そこで「構造体」の出番です。関連する変数を1つの型としてまとめられます。

\begin{goodexample}{構造体でまとめる}
\begin{lstlisting}
struct LedData {
    bool state = false;
    uint8_t brightness = 0;
};

LedData led1;  // LED1のデータ一式
LedData led2;  // LED2のデータ一式

led1.state = true;       // ドット(.)でアクセス
led2.brightness = 128;
\end{lstlisting}
\end{goodexample}

構造体を使えば、「LEDに必要なデータは何か」が型として明確になり、LEDが何個増えても変数の定義は1行で済みます。

\begin{cautionbox}[構造体の定義だけでは「実体」はない]
初学者がつまずきやすいポイントです。\texttt{struct LedData \{...\};} と書いた時点では、メモリ上に何も作られていません。これは「設計図（型）」を定義しただけです。

\begin{lstlisting}
// これは「設計図」を書いただけ（実体はまだない）
struct LedData {
    bool state = false;
    uint8_t brightness = 0;
};

// ここで初めて「実体（インスタンス）」が作られる
LedData led1;   // LedData型の実体を1つ作成
LedData led2;   // もう1つ作成（led1とled2は別々の実体）
\end{lstlisting}

たとえるなら、構造体の定義は「たい焼きの型（金型）」です。型を作っただけではたい焼きは食べられません。型に生地を流し込んで焼いて初めて、食べられるたい焼き（実体）ができます。

\begin{itemize}
    \item \textbf{型（設計図）}: \texttt{struct LedData \{...\};} — 「こういうデータの組み合わせがある」と定義
    \item \textbf{実体（インスタンス）}: \texttt{LedData led1;} — 実際にメモリ上に作られた「もの」
\end{itemize}

1つの型から何個でも実体を作れます。led1とled2はどちらもLedData型ですが、それぞれ独立した別のデータを持ちます。
\end{cautionbox}

\end{qabox}

%% ── Q2: クラス ──────────────────────────────
\begin{qabox}{構造体とクラスって何が違うの？}

構造体は「データをまとめる」ものでした。クラスはさらに「そのデータを操作する関数」も一緒にまとめます。

\begin{goodexample}{クラス = データ + 関数}
\begin{lstlisting}
class LedModule {
private:
    uint8_t _pin;       // データ（メンバ変数）
    bool _state;

public:
    // コンストラクタ（後述）
    LedModule(uint8_t pin) : _pin(pin), _state(false) {}

    // メンバ関数
    void turnOn() {
        _state = true;
        digitalWrite(_pin, HIGH);
    }

    void turnOff() {
        _state = false;
        digitalWrite(_pin, LOW);
    }
};
\end{lstlisting}
\end{goodexample}

構造体でもメンバ関数を持てますが、C++の慣習として次のように使い分けます。

\begin{itemize}
  \item \textbf{struct} — データの入れ物（Config、Dataなど）
  \item \textbf{class} — データ＋操作をセットにしたもの（モジュールなど）
\end{itemize}

\begin{notebox}[structとclassの技術的な違い]
C++ではstructとclassの唯一の違いは「デフォルトのアクセス修飾子」です。structはデフォルトpublic、classはデフォルトprivateです。本書ではデータの入れ物にstruct、操作を持つものにclassを使います。
\end{notebox}

\begin{cautionbox}[クラスも「型を定義している」だけ]
構造体と同じく、クラスの定義だけでは実体は作られません。

\begin{lstlisting}
// これは設計図（型）
class LedModule {
    uint8_t _pin;
    bool _state;
public:
    LedModule(uint8_t pin) : _pin(pin), _state(false) {}
    void turnOn();
    void turnOff();
};

// ここで実体（インスタンス）が作られる
LedModule led1(13);   // ピン13のLEDモジュール
LedModule led2(14);   // ピン14のLEDモジュール（別の実体）
\end{lstlisting}

\texttt{led1.turnOn()} と呼ぶと \texttt{led1} だけが点灯し、\texttt{led2} には影響しません。同じ設計図から作られていても、実体はそれぞれ独立しています。

この「型の定義」と「実体の作成」が別であるという点は、本書の設計パターン全体を通じて重要な考え方です。Config構造体の「型定義」とConfigインスタンスの「値定義」が別ファイルになっているのも、この原則に基づいています。
\end{cautionbox}

\end{qabox}

%% ── Q3: コンストラクタ ──────────────────────
\begin{qabox}{コンストラクタって何？}

コンストラクタは「オブジェクトが作られるときに自動で呼ばれる関数」です。初期化忘れを防げます。

\begin{badexample}{初期化忘れの危険}
\begin{lstlisting}
LedModule led;         // ピン番号が未設定！
led.turnOn();          // どのピンが光る？→バグ
\end{lstlisting}
\end{badexample}

\begin{goodexample}{コンストラクタで確実に初期化}
\begin{lstlisting}
class LedModule {
private:
    uint8_t _pin;
    bool _state;

public:
    // コンストラクタ — 初期化子リストで初期化
    LedModule(uint8_t pin) : _pin(pin), _state(false) {
        // ピンを指定しないと作れない
    }
};

LedModule led(13);   // OK: ピン13で作成
// LedModule led2;   // コンパイルエラー！ピン未指定
\end{lstlisting}
\end{goodexample}

コロン（\texttt{:}）の後に書く \texttt{\_pin(pin), \_state(false)} の部分を「初期化子リスト」と呼びます。メンバ変数を効率よく初期化する書き方です。
\end{qabox}

%% ── Q4: private/public ──────────────────────
\begin{qabox}{private/publicって何？}

クラスの中身を「外から触れる部分」と「触れない部分」に分けるしくみです。

\vspace{2mm}
イメージとしては、銀行の窓口を思い浮かべてください。

\begin{itemize}
  \item \textbf{public}（窓口） — お客さん（外部のコード）がアクセスできる
  \item \textbf{private}（金庫） — 銀行員（クラス内部）だけがアクセスできる
\end{itemize}

\begin{lstlisting}
class LedModule {
private:
    uint8_t _pin;    // 外から直接変更させない

public:
    LedModule(uint8_t pin) : _pin(pin) {}
    void turnOn();   // 外から呼べる
};

LedModule led(13);
led.turnOn();      // OK: publicなので呼べる
// led._pin = 99;  // コンパイルエラー！privateには触れない
\end{lstlisting}

なぜわざわざ隠すのか？ 外部から \texttt{\_pin} を勝手に変更されると、LEDモジュールの動作が壊れる可能性があるからです。「触ってほしくないものはprivateにする」のが安全設計の基本です。

\begin{notebox}[命名の慣習]
privateメンバ変数には先頭にアンダースコア（\texttt{\_}）を付ける慣習があります。\texttt{\_pin}、\texttt{\_state}のように書くと、「これは内部変数だな」と一目で分かります。
\end{notebox}

\end{qabox}

%% ── Q5: virtualの予告 ──────────────────────
\begin{qabox}{virtualって何？（予告）}

\texttt{virtual}は「クラスの継承」という仕組みで使うキーワードです。これを使うと、異なるクラスを同じ型として扱えるようになります。

……と言われてもピンと来ないと思います。

大丈夫です。\texttt{virtual}の詳しい説明は本編の「インターフェースで統一する」の章で、具体例を交えて解説します。ここでは「そういうキーワードがあるんだな」とだけ覚えておいてください。

\begin{notebox}[ここまでのまとめ]
\begin{itemize}
  \item \textbf{構造体} — 関連する変数をまとめる箱
  \item \textbf{クラス} — データ＋操作をセットにしたもの
  \item \textbf{コンストラクタ} — オブジェクト生成時に自動実行される初期化関数
  \item \textbf{private/public} — アクセス範囲を制限して安全性を確保
  \item \textbf{virtual} — 「インターフェースで統一する」の章で詳しく！
\end{itemize}
\end{notebox}

これだけ分かっていれば、本編を読み進められます。さあ、次の章に進みましょう！

\end{qabox}
